% TEMPLATE-MILC-v3.tex - plantilla maestra para todas las guias MILC v3
% Ver scripts/generadores/build-guia-milc.py para la lista completa de
% claves esperadas. El script valida que no queden placeholders sin
% reemplazar antes de escribir el archivo final.

\documentclass[11pt,letterpaper]{article}
\usepackage[margin=1.75cm]{geometry}
\usepackage[table]{xcolor}
\usepackage{fontspec}
\usepackage[spanish]{babel}
\usepackage{tikz}
\usepackage[most]{tcolorbox}
\usepackage{tabularx}
\usepackage{array}
\usepackage{fancyhdr}
\usepackage{titlesec}
\usepackage{ragged2e}
\usepackage{enumitem}
\usepackage{microtype}

\setmainfont{Helvetica Neue}
\setsansfont{Helvetica Neue}
\setlength{\parindent}{0pt}
\setlength{\parskip}{6pt}
\hyphenpenalty=200
\exhyphenpenalty=200
\pretolerance=400
\tolerance=2000
\setlength{\emergencystretch}{1em}
\renewcommand{\arraystretch}{1.25}
\setlist{leftmargin=5mm,itemsep=1mm,topsep=1mm}
\newcolumntype{Y}{>{\RaggedRight\arraybackslash}X}

\definecolor{milcMagenta}{HTML}{E6007E}
\definecolor{milcVino}{HTML}{5A0038}
\definecolor{milcTurquesa}{HTML}{0093A5}
\definecolor{milcVerde}{HTML}{58A923}
\definecolor{milcMostaza}{HTML}{E5B400}
\definecolor{milcOcre}{HTML}{B86600}
\definecolor{milcNegro}{HTML}{241816}
\definecolor{milcGris}{HTML}{F7F7F4}
\definecolor{milcLinea}{HTML}{E8E2DD}
\definecolor{milcGradePrimary}{HTML}{0066FF}
\definecolor{milcGradeDark}{HTML}{003D99}
\definecolor{milcGradeSoft}{HTML}{D6E8FF}
\definecolor{milcGradeAccent}{HTML}{CFE7FF}
\definecolor{milcGradeText}{HTML}{FFFFFF}
\color{milcNegro}

\pagestyle{fancy}
\fancyhf{}
\lhead{\small Tecnología e Informática · Grado 6}
\rhead{\small Guía 6 · MILC v3}
\cfoot{\small\thepage}
\renewcommand{\headrulewidth}{0pt}

\newtcolorbox{titlebox}[1]{
  enhanced,colback=#1,colframe=#1,arc=7mm,boxrule=0pt,
  left=7mm,right=7mm,top=3mm,bottom=3mm,
  before skip=8pt,after skip=8pt,
  fontupper=\sffamily\bfseries\Large\color{white}
}

\newtcolorbox{softbox}[3]{
  enhanced,colback=#2,colframe=#1,borderline west={4pt}{0pt}{#1},
  arc=2mm,boxrule=.4pt,left=4mm,right=4mm,top=3mm,bottom=3mm,
  before skip=6pt,after skip=8pt,title={#3},coltitle=white,
  colbacktitle=#1,fonttitle=\sffamily\bfseries,fontupper=\RaggedRight,
  attach boxed title to top left={xshift=3mm,yshift=-2mm},
  boxed title style={arc=2mm, boxrule=0pt}
}

\newtcolorbox{drawbox}[2]{
  enhanced,colback=white,colframe=#1,arc=4mm,boxrule=1pt,
  left=5mm,right=5mm,top=4mm,bottom=4mm,before skip=8pt,after skip=8pt,
  title={#2},coltitle=white,colbacktitle=#1,fonttitle=\sffamily\bfseries,
  fontupper=\RaggedRight,
  attach boxed title to top left={xshift=4mm,yshift=-2mm},
  boxed title style={arc=3mm, boxrule=0pt}
}

\newtcolorbox{infoband}[1]{
  enhanced,colback=#1!8,colframe=#1!50,arc=2mm,boxrule=.5pt,
  left=4mm,right=4mm,top=2mm,bottom=2mm,before skip=4pt,after skip=4pt,
  fontupper=\small\RaggedRight
}

% Puente narrativo entre secciones (contrato editorial Conéctate).
% Da continuidad: "Acabas de X. Ahora vas a Y porque Z."
% Visualmente: banda izquierda en color del grado, texto en itálica pequeña.
\newtcolorbox{puentebox}{
  enhanced,colback=milcGradeSoft,colframe=milcGradeDark,
  borderline west={3pt}{0pt}{milcGradeDark},
  arc=0mm,boxrule=0pt,left=5mm,right=4mm,top=2.5mm,bottom=2.5mm,
  before skip=4pt,after skip=8pt,
  fontupper=\itshape\small\color{milcNegro}\RaggedRight
}
\newcommand{\puente}[1]{%
  \begin{puentebox}%
    \textcolor{milcGradeDark}{\textbf{\textrightarrow}}\hspace{2mm}#1%
  \end{puentebox}%
}

\newcommand{\linead}[1]{\noindent\color{milcLinea}\rule{#1}{0.5pt}\color{milcNegro}}
\newcommand{\respuesta}[1]{\par\vspace{2mm}\foreach \n in {1,...,#1}{\noindent\color{milcLinea}\rule{\linewidth}{0.5pt}\par\vspace{5mm}}\color{milcNegro}}
\newcommand{\checkbox}{\textcolor{milcGradeDark}{\fbox{\phantom{x}}}\hspace{5pt}}

\newcommand{\phasecard}[4]{
\begin{tcolorbox}[enhanced,colback=#3,colframe=#2,arc=3mm,boxrule=.5pt,height=3.05cm,valign=center,halign=center,left=1.5mm,right=1.5mm,top=2mm,bottom=2mm]
{\sffamily\bfseries\fontsize{22}{24}\selectfont\textcolor{#2}{#1}}\par
{\sffamily\bfseries\small #4}
\end{tcolorbox}
}

\begin{document}
\thispagestyle{empty}
\begin{tikzpicture}[remember picture,overlay]
  \fill[white] (current page.south west) rectangle (current page.north east);
  \path[fill=milcGradePrimary,rounded corners=16mm]
    ([xshift=.55cm,yshift=-.55cm]current page.north west) rectangle
    ([xshift=-.55cm,yshift=3.65cm]current page.south east);

  \node[anchor=north west,text=milcGradeText] at ([xshift=2.0cm,yshift=-1.85cm]current page.north west)
    {\sffamily\bfseries\fontsize{14}{16}\selectfont GRADO};
  \node[anchor=north west,text=milcGradeText,opacity=.82] at ([xshift=2.0cm,yshift=-2.75cm]current page.north west)
    {\sffamily\bfseries\fontsize{9}{11}\selectfont SERIE GUÍAS · TECNOLOGÍA E INFORMÁTICA};

  \begin{scope}[shift={([xshift=-3.05cm,yshift=-2.55cm]current page.north east)}]
    \fill[white,opacity=.80] (-.62,-.52) rectangle (.62,.52);
    \draw[milcGradeText,line width=1pt] (-.62,-.52) rectangle (.62,.52);
    \fill[milcGradeDark] (-.42,-.38) rectangle (-.22,.06);
    \fill[milcGradeAccent] (-.10,-.38) rectangle (.10,.24);
    \fill[milcGradePrimary!60!black] (.22,-.38) rectangle (.42,.38);
  \end{scope}

  \node[anchor=west,text=milcGradeText] at ([xshift=1.9cm,yshift=-12.85cm]current page.north west)
    {\sffamily\bfseries\fontsize{124}{124}\selectfont 6};
  \draw[milcGradeText,line width=1.7pt]
    ([xshift=4.52cm,yshift=-11.68cm]current page.north west) circle (.13cm);

  \node[anchor=west,text=milcGradeText,text width=8.7cm,align=left] at ([xshift=10.35cm,yshift=-11.6cm]current page.north west)
    {
     {\sffamily\bfseries\fontsize{19}{22}\selectfont Mi huella\\digital ---\\el rastro\\que dejas\\sin notarlo}\\[4mm]
     {\sffamily\bfseries\fontsize{23}{26}\selectfont Guía 6-6-TIC}\\[2mm]
     {\sffamily\fontsize{13}{16}\selectfont Período 1 · Comunicación e identidad digital}\\[5mm]
     {\sffamily\bfseries\fontsize{11}{13}\selectfont Institución Educativa Sor María Juliana}\\[1mm]
     {\sffamily\fontsize{10}{12}\selectfont Autor: Dr. Álvaro Cárdenas Orozco}
    };

  \path[fill=milcGradeSoft,rounded corners=7mm]
    ([xshift=1.35cm,yshift=.85cm]current page.south west) rectangle
    ([xshift=-1.35cm,yshift=4.25cm]current page.south east);
  \draw[milcGradeDark,line width=.65pt,rounded corners=7mm]
    ([xshift=1.35cm,yshift=.85cm]current page.south west) rectangle
    ([xshift=-1.35cm,yshift=4.25cm]current page.south east);
  \node[anchor=center,text=milcNegro,text width=17.0cm,align=center] at ([yshift=2.55cm]current page.south)
    {\sffamily\fontsize{10}{12}\selectfont Metodología MILC · Apertura ancestral · Pensamiento computacional · Triángulo Dussel-Estoicismo-Floridi\\
    \textcolor{milcGradeDark}{\bfseries Producto:} Un \textbf{mapa personal de la huella digital de un día tuyo}: en una hoja del cuaderno dibujas un \emph{reloj de 24 horas} en círculo, y en cada hora marcas \textbf{qué apps abriste, qué buscaste, dónde estuviste} (según tu celular). Después clasificas cada huella en 3 colores: \textbf{verde} (huella que tú decidiste dejar), \textbf{amarillo} (huella que dejaste sin pensar), \textbf{rojo} (huella que tal vez te conviene borrar o evitar). Más \textbf{3 compromisos escritos} para reducir las huellas rojas en la próxima semana.};
\end{tikzpicture}
\mbox{}
\newpage

\section*{Apertura: Mi huella digital --- el rastro que dejas sin darte cuenta}

\begin{softbox}{milcGradeDark}{milcGradeSoft}{Saber ancestral}
\textbf{En el río Cauca el pescador sabía dónde estaban los peces porque leía huellas.} Antes de los sonares y los celulares, el pescador del Pacífico y del Cauca leía el agua para saber qué había debajo: las \textbf{ondas pequeñas} mostraban un cardumen, las \textbf{burbujas} contaban dónde dormía el bocachico, las \textbf{ramas mordidas} de la orilla decían \emph{``aquí estuvo un capibara hace dos horas''}. Cada criatura, sin querer, dejaba un rastro. Por ese rastro la gente sabía exactamente quién había pasado, cuándo y a dónde iba. En internet pasa parecido: \textbf{cada vez que entras a una app, buscas algo, das me gusta, dejas un rastro}. No lo ves, pero está. Otros sí lo leen --- algunos para cuidarte (tus papás), otros para venderte cosas (las empresas), y a veces personas malas para hacerte daño.
\end{softbox}

\begin{softbox}{milcTurquesa}{milcGris}{Saber contemporáneo}
Tu \textbf{huella digital} es el rastro que dejas cuando usas internet. \textbf{No es lo mismo que la identidad digital.} La identidad es lo que TÚ decides mostrar de ti (tu foto, tu apodo, tus videos). La huella es \textbf{todo lo demás}: lo que buscas, dónde estuviste según el GPS, cuánto tiempo viste cada TikTok, qué apps usas, en qué horarios, hasta cómo deslizas el dedo. Las empresas guardan esa huella para venderte cosas. \textbf{Tú casi no la ves}, pero ahí queda. Hay 3 tipos de huella, y hoy las vas a reconocer. Cuando entiendes tu huella, puedes \emph{borrar las que no te sirven, cuidar las que sí}, y empezar a moverte por internet con criterio en vez de a ciegas.
\end{softbox}

\begin{softbox}{milcMostaza}{milcGris}{Pregunta-puente}
\textit{Si tu mamá pudiera ver \textbf{todo lo que hiciste en el celular ayer} --- cada video, cada búsqueda, cada chat, cada lugar donde estuviste según el GPS --- ¿le gustaría lo que vería? ¿O hay cosas que preferirías que no viera?}
\end{softbox}

\begin{softbox}{milcVerde}{milcGris}{Saber y saber hacer}
Al terminar la clase: (1) podrás \textbf{identificar} 3 tipos distintos de huella digital; (2) sabrás \textbf{explicar} la diferencia entre identidad digital y huella digital; (3) habrás \textbf{evaluado} las huellas de un día tuyo y las clasificarás en verde-amarillo-rojo; (4) habrás \textbf{escrito} 3 compromisos para reducir tus huellas rojas.
\end{softbox}

\small
\begin{tabularx}{\textwidth}{>{\bfseries}p{3.0cm}Y}
\rowcolor{milcGradePrimary!12}Autor & Dr. Álvaro Cárdenas Orozco\\
DBA & Reconoce principios y conceptos propios de la tecnología, relacionando artefactos con su utilización segura (MEN, T\&I 6°)\\
Referentes & Saber ancestral del pregonero · Pensamiento computacional inicial · Cuidado de la palabra\\
Producto final & Un \textbf{mapa personal de la huella digital de un día tuyo}: en una hoja del cuaderno dibujas un \emph{reloj de 24 horas} en círculo, y en cada hora marcas \textbf{qué apps abriste, qué buscaste, dónde estuviste} (según tu celular). Después clasificas cada huella en 3 colores: \textbf{verde} (huella que tú decidiste dejar), \textbf{amarillo} (huella que dejaste sin pensar), \textbf{rojo} (huella que tal vez te conviene borrar o evitar). Más \textbf{3 compromisos escritos} para reducir las huellas rojas en la próxima semana.\\
\end{tabularx}
\normalsize

\puente{Hoy haces 4 cosas: identificas las apps que usaste ayer, descubres los 3 tipos de huella, evalúas tus huellas, y firmas compromisos.}

\begin{titlebox}{milcGradeDark}
Ruta MILC: así vamos a aprender
\end{titlebox}
\begin{tabularx}{\textwidth}{>{\centering\arraybackslash}X>{\centering\arraybackslash}X>{\centering\arraybackslash}X>{\centering\arraybackslash}X}
\phasecard{1}{milcVerde}{milcGris}{Escuta\\[-1mm]\small Observo, escucho y pregunto.} &
\phasecard{2}{milcTurquesa}{milcGris}{Sistematización\\[-1mm]\small Pensamiento computacional.} &
\phasecard{3}{milcMagenta}{milcGris}{Praxis\\[-1mm]\small Construyo y aplico.} &
\phasecard{4}{milcVino}{milcGris}{Evaluación\\[-1mm]\small Reflexión liberadora.}
\end{tabularx}


\puente{Empezamos con un inventario de tu día de ayer en el celular.}

\begin{titlebox}{milcVerde}
Fase 1 · Escuta
\end{titlebox}

\begin{softbox}{milcVerde}{milcGris}{Escena para escuchar}
\textbf{Actividad 1 · IDENTIFICA --- Inventario de tu día de ayer} (12 min · individual). En el cuaderno, dibuja una tabla de 2 columnas. En la columna izquierda escribe estas \textbf{6 horas} de un día normal: \emph{6 am (despertar), 8 am (camino al cole), 12 m (descanso), 4 pm (al volver a casa), 7 pm (después de comer), 10 pm (antes de dormir)}. En la columna derecha escribe \textbf{qué hiciste en el celular en cada hora}: ¿qué app abriste? ¿qué buscaste? ¿con quién chateaste? ¿qué video viste? Sé honesto; este cuaderno no se publica. Si no usaste el celular en una hora, escribe \emph{``no lo toqué''} --- esa también es una respuesta.
\end{softbox}

\checkbox Hice la tabla de 6 horas.\\
\checkbox Anoté con honestidad qué hice en el celular en cada una.\\
\checkbox Si en alguna hora no usé el celular, lo anoté también.

\begin{infoband}{milcVerde}
\textbf{Lo que va al cuaderno (Actividad 1 · IDENTIFICA).} Título: «Actividad 1 --- Mi día de ayer en el celular». \textbf{Formato:} tabla de 2 columnas, 6 filas. \textbf{Extensión:} 6 filas con honestidad. \textbf{Sabes que terminaste cuando} tu tabla refleja un día real tuyo, no un día imaginario \emph{``perfecto''}.
\end{infoband}

\puente{Después del inventario, aprendes los 3 tipos de huella digital.}

\begin{titlebox}{milcTurquesa}
Fase 2 · Sistematización con pensamiento computacional
\end{titlebox}

\begin{softbox}{milcTurquesa}{milcGris}{Cuatro pilares aplicados al tema}
Lo que acabaste de hacer es \textbf{ver tu huella}. Casi nunca te detienes a verla. Ahora vas a aprender que esa huella tiene 3 tipos distintos. Saberlos te ayuda a saber cuál cuidar y cuál borrar.
\end{softbox}

\small
\begin{tabularx}{\textwidth}{>{\bfseries\RaggedRight\arraybackslash}p{3.6cm}Y}
\rowcolor{milcTurquesa!90}\color{white}Pilar & \color{white}\bfseries Aplicación al tema\\
1. Descomposición & \textbf{Tipo 1 --- Huella activa.} Es la huella que dejas \emph{a propósito}, porque tú decidiste dejarla. \textbf{Ejemplos:} publicar un TikTok, subir una historia a Instagram, escribir un comentario, mandar un mensaje en WhatsApp. \textbf{Lo importante:} esta huella la controlas tú. Si eres cuidadoso al publicar, no hay problema. Si publicas sin pensar, te queda.\\
2. Reconocimiento de patrones & \textbf{Tipo 2 --- Huella pasiva.} Es la huella que dejas \emph{sin darte cuenta}, simplemente por usar las apps. \textbf{Ejemplos:} TikTok sabe cuántos segundos ves cada video, Google guarda cada cosa que buscas, Maps sabe dónde estuviste, Instagram registra cuántas veces abriste su app y a qué hora. \textbf{Lo importante:} esta huella casi no la ves, pero las empresas sí. Sirve para venderte cosas.\\
3. Abstracción & \textbf{Tipo 3 --- Huella secundaria.} Es la huella que dejan \emph{otros sobre ti}. \textbf{Ejemplos:} un amigo te toma una foto y la sube; alguien te menciona en un comentario; tu prima te etiqueta en una historia; un familiar publica fotos tuyas de pequeño. \textbf{Lo importante:} esta huella no la controlas directamente. Pero puedes pedirle a la persona que la baje, o usar opciones de privacidad para que no te etiqueten sin permiso.\\
4. Algoritmo conceptual & \textbf{Las 4 reglas para cuidar tu huella.} (1) \textbf{Revisa los permisos:} entra a Ajustes del celular y mira qué apps tienen acceso al GPS, micrófono y cámara. Quita los que no tengan sentido (¿por qué una linterna necesita GPS?). (2) \textbf{Borra de vez en cuando:} el historial de búsqueda, los mensajes viejos, las cuentas que ya no usas. Lo que no borras, se acumula. (3) \textbf{Piensa antes de publicar (huella activa):} si tu mamá no quisiera verlo, no lo publiques. (4) \textbf{Pide permiso, da permiso:} antes de publicar fotos de otros, pides; antes de aceptar que te etiqueten, miras dónde te etiquetan.\\
\end{tabularx}
\normalsize

\begin{drawbox}{milcTurquesa}{Los 3 tipos de huella digital}
\textbf{1. Huella activa} --- la que dejas a propósito (publicar, comentar). \\[1mm]
\textbf{2. Huella pasiva} --- la que dejas sin notar (búsquedas, ubicación, tiempo en apps). \\[1mm]
\textbf{3. Huella secundaria} --- la que dejan otros sobre ti (fotos donde te etiquetan, comentarios).
\end{drawbox}

\begin{softbox}{milcMostaza}{milcGris}{Errores comunes y su corrección}
\textbf{Mitos sobre la huella digital que escucharás.} (1) \emph{``Si uso modo incógnito, nadie me ve''} --- Falso. El modo incógnito esconde el historial \emph{de tu computador}, pero la app y el sitio web te siguen viendo igual. (2) \emph{``Si borro la app, se borra todo''} --- Falso. La empresa guarda tus datos aunque tú borres la app. Tienes que pedir formalmente que borren tu cuenta. (3) \emph{``Yo no soy famoso, a nadie le importa mi huella''} --- Falso. A las empresas SÍ les importa: tus datos se venden a publicistas. Tu huella vale dinero. (4) \emph{``Si no tengo nada que esconder, no importa''} --- Falso. Privacidad no es esconder; es escoger qué muestras y a quién.
\end{softbox}

\begin{infoband}{milcTurquesa}
\textbf{Lo que va al cuaderno (Actividad 2 · EXPLICA).} Título: «Actividad 2 --- Los 3 tipos de huella con ejemplos míos». \textbf{Formato:} tabla de 2 columnas, 3 filas. Columna 1: nombre del tipo. Columna 2: \emph{un ejemplo tuyo} de huella de ese tipo del día de ayer (saca del inventario que ya hiciste). \textbf{Extensión:} 3 filas. \textbf{Sabes que terminaste cuando} podrías explicarle a un primo más pequeño qué es huella activa, pasiva y secundaria.
\end{infoband}

\puente{Con los 3 tipos claros, dibujas tu reloj de 24 horas y clasificas tus huellas por color.}

\begin{titlebox}{milcMagenta}
Fase 3 · Praxis con pensamiento computacional
\end{titlebox}

\begin{softbox}{milcMagenta}{milcGris}{Construyendo el producto con los cuatro pilares}
\textbf{Actividad 3 · EVALÚA --- Tu reloj de 24 horas con colores} (25 min · individual). Vas a dibujar en una hoja del cuaderno tu \textbf{huella digital de un día} en forma de reloj. Lo dividirás en 24 horas, marcarás qué hiciste en cada una con el celular, y \textbf{clasificarás cada acción con un color}: verde (sano), amarillo (cuidado), rojo (peligro). Al final firmas 3 compromisos.
\end{softbox}

\small
\begin{tabularx}{\textwidth}{>{\bfseries\RaggedRight\arraybackslash}p{3.6cm}Y}
\rowcolor{milcMagenta!90}\color{white}Pilar & \color{white}\bfseries Aplicación al producto\\
Descomposición & \textbf{Paso 1.} Dibuja un \textbf{círculo grande en el centro de la hoja}. Divídelo en 24 partes (como un reloj pero con 24 horas en vez de 12). Marca las horas en el borde: 6am, 8am, 10am, 12m, 2pm, 4pm, 6pm, 8pm, 10pm, 12am, etc.\\
Patrones & \textbf{Paso 2.} En cada hora del círculo, escribe (o dibuja con un ícono) \textbf{lo que hiciste con el celular} en esa hora ayer. Usa la tabla de la Actividad 1 como base. Si no usaste el celular en una hora, deja la hora en blanco --- ese vacío también cuenta.\\
Abstracción & \textbf{Paso 3.} Ahora viene la parte importante. Pinta \textbf{el borde de cada hora con un color} según esta regla: \textbf{verde} si lo que hiciste fue sano (chat con familia, estudiar, jugar con amigos conocidos). \textbf{Amarillo} si fue dudoso (ver muchos videos seguidos, perder horas sin notar, chats con desconocidos). \textbf{Rojo} si fue riesgoso (compartir info personal, hablar con extraños, ver cosas que no deberías).\\
Algoritmo de armado & \textbf{Paso 4.} Cuenta los colores. ¿Cuántas horas verdes hay? ¿Cuántas amarillas? ¿Cuántas rojas? Anótalo al lado del círculo: \emph{``Verde: 4h. Amarillo: 6h. Rojo: 1h.''}. Sin trampa --- cuenta lo que realmente hiciste.\\
\end{tabularx}
\normalsize

\begin{drawbox}{milcMagenta}{Antes de mostrar el mapa}
\checkbox El reloj de 24 horas ocupa el centro de la hoja. \\[1mm]
\checkbox Cada hora tiene anotado lo que hice (o un vacío si no usé el celular). \\[1mm]
\checkbox Cada hora está pintada con verde, amarillo o rojo. \\[1mm]
\checkbox Conté cuántas horas tuve de cada color. \\[1mm]
\checkbox Identifiqué 3 acciones a reducir (de amarillas o rojas). \\[1mm]
\checkbox Escribí 3 compromisos concretos y los firmé.
\end{drawbox}

\begin{softbox}{milcMagenta}{milcGris}{Plantilla de trabajo}
\textbf{Estructura del reloj:} círculo grande, 24 divisiones, hora marcada en el borde, acción adentro, color en el borde. \textbf{Estructura del compromiso:} \emph{``Acción a reducir [N]: [qué hago hoy]. Compromiso: [qué voy a hacer en cambio, medible].''}.
\end{softbox}

\begin{infoband}{milcMagenta}
\textbf{Lo que va al cuaderno (Actividad 3 · EVALÚA).} Título: «Actividad 3 --- Mi huella digital de un día». \textbf{Formato:} reloj circular de 24 horas con colores + lista de 3 compromisos. \textbf{Extensión:} 1 hoja entera del cuaderno. \textbf{Sabes que terminaste cuando} mirando el reloj puedes decir en voz alta \emph{``estas son las horas que necesito cuidar más''} y tienes 3 compromisos medibles para hacerlo.
\end{infoband}

\puente{El producto es ese reloj clasificado más los 3 compromisos para reducir huellas rojas.}

\begin{titlebox}{milcMostaza}
Producto de la sesión
\end{titlebox}
\textbf{Reloj de 24 horas + 3 compromisos firmados}.
\begin{softbox}{milcMostaza}{milcGris}{Criterios de calidad}
(1) El reloj cubre las 24 horas del día. (2) Cada hora con actividad tiene su anotación. (3) Las horas están clasificadas en verde, amarillo o rojo. (4) El conteo de colores está hecho con honestidad. (5) Los 3 compromisos son medibles (con tiempo, con número, con acción concreta).
\end{softbox}

\puente{Antes de salir, miras tus 3 compromisos y verificas que sean concretos (no \emph{``voy a portarme mejor''}, sino \emph{``voy a desactivar el GPS de TikTok''}).}

\begin{titlebox}{milcVino}
Fase 4 · Evaluación liberadora — cinco dimensiones
\end{titlebox}

\begin{softbox}{milcVino}{milcGris}{Sentido de la evaluación}
Evaluar no es solo revisar una nota. Es reconocer qué aprendí, qué cambió en mí, qué emociones manejé, cómo me relaciono con la ciudad y la comunidad, y qué le dejaré a quien viene detrás.
\end{softbox}

\textbf{Mi reflexión liberadora en cinco dimensiones.} Completa cada frase iniciada:

\small
\begin{tabularx}{\textwidth}{>{\bfseries\RaggedRight\arraybackslash}p{3.4cm}Y>{\RaggedRight\arraybackslash}p{4.6cm}}
\rowcolor{milcVino}\color{white}Dimensión & \color{white}\bfseries Mi respuesta (frame starter) & \color{white}\bfseries Algo que descubrí\\
1. Personal & Lo que más cambió en mí fue \linead{6.5cm} & \linead{4.2cm}\\
2. Emocional & La emoción más fuerte fue \linead{2.5cm} y la regulé \linead{3cm} & \linead{4.2cm}\\
3. Ciudadana & En democracia esto importa porque \linead{6cm} & \linead{4.2cm}\\
4. Local (Cartago) & En mi barrio sería útil para \linead{6.5cm} & \linead{4.2cm}\\
5. Intergeneracional & A mi abuela/abuelo le diría \linead{6.5cm} & \linead{4.2cm}\\
\end{tabularx}
\normalsize

\begin{infoband}{milcVino}
\textbf{Para la próxima sustentación.} Vuelve a esta tabla en seis meses. Verás que las respuestas cambian — no porque lo aprendido cambie, sino porque tú cambias. La evaluación liberadora es una conversación contigo a través del tiempo.
\end{infoband}


\puente{Cerramos con 3 voces que te ayudan a entender por qué te conviene saber qué rastro dejas.}

\begin{titlebox}{milcGradeDark}
Triángulo de pensamiento · cierre filosófico
\end{titlebox}

\begin{softbox}{milcVino}{milcGris}{Dussel · lente del nosotros}
\textit{````No eres lo que las empresas creen que eres por tus búsquedas. Eres más que tu huella.''''}

Las empresas que recogen tu huella creen que pueden predecir todo de ti: qué te gusta, qué vas a comprar, qué te asusta. \textbf{Pero tú eres mucho más} que esa huella. La huella es una sombra muy pequeña de quien eres. Conocer tu huella te ayuda a no dejar que esa sombra te defina. Tú decides qué partes del río de tu vida quieres que se vean.

\textbf{¿Estoy \textbf{dejándome describir} por mis búsquedas y mis likes, o sigo siendo yo el que decide quién soy?}
\end{softbox}
\linead{\linewidth}\\[1mm]
\linead{\linewidth}

\begin{softbox}{milcMostaza}{milcGris}{Estoicismo · Marco Aurelio · cuidado interior}
\textit{````Lo que haces a solas habla más fuerte de ti que lo que muestras en público.''''}

Marco Aurelio escribió esto hace 1800 años, pero parece escrito ayer. \textbf{Tu huella digital es lo que haces ``a solas'' con el celular}: nadie te ve, pero el celular sí. Cuida tu huella como cuidarías tu cuaderno privado: si tu mamá lo abriera, ¿te sentirías bien con lo que ve? Esa pregunta es el espejo.

\textbf{¿Hago en el celular cosas que coinciden con quien soy, o cosas que me darían pena si las viera mi familia?}
\end{softbox}
\linead{\linewidth}\\[1mm]
\linead{\linewidth}

\begin{softbox}{milcTurquesa}{milcGris}{Floridi · lente de la infoesfera}
\textit{````Privacidad no es esconder; privacidad es decidir qué se cuenta de ti y a quién.''''}

Mucha gente dice \emph{``no tengo nada que esconder''} para justificar no cuidar la huella. Pero \textbf{privacidad no es esconder}: privacidad es \textbf{tú escoges} qué muestras, a quién, y cuándo. Es libertad. Cuando alguien recoge tu huella sin que tú decidas, te están quitando esa libertad. Cuidar la huella es defender la libertad de seguir decidiendo.

\textbf{Cuando entrego mis datos sin pensar, ¿estoy regalando libertad sin notarlo?}
\end{softbox}
\linead{\linewidth}\\[1mm]
\linead{\linewidth}

\begin{titlebox}{milcVino}
Compromiso final
\end{titlebox}

\small
\begin{tabularx}{\textwidth}{>{\RaggedRight\arraybackslash}p{3.0cm}YYY}
\rowcolor{milcVino}\color{white}\bfseries Criterio & \color{white}\bfseries Lo logré & \color{white}\bfseries En proceso & \color{white}\bfseries Necesito apoyo\\
Comprensión del tema & Explico ideas centrales con mis palabras. & Reconozco ideas pero falta precisión. & Necesito repasar conceptos base.\\
Pensamiento computacional & Apliqué los 4 pilares al diseñar mi producto. & Apliqué algunos pilares. & Necesito releer la fase 2.\\
Aplicación y producto & Mi producto cumple los criterios y se puede revisar. & Producto funcional con ajustes pendientes. & Aún en construcción.\\
Reflexión 5 dimensiones & Respondí cada dimensión con honestidad. & Respondí algunas con detalle. & Mis respuestas son muy generales.\\
Triángulo de pensamiento & Conecté las 3 voces con mi trabajo real. & Conecté al menos una. & No relaciono las voces con mi proceso.\\
\end{tabularx}
\normalsize

\textbf{Hoy entendí que:}\\[1mm]
\linead{\linewidth}\\[1mm]
\linead{\linewidth}

\textbf{Mi compromiso para la próxima sesión es:}\\[1mm]
\linead{\linewidth}\\[1mm]
\linead{\linewidth}

\begin{softbox}{milcMostaza}{milcGris}{Cita de despedida · Marco Aurelio}
\textit{``El obstáculo en el camino se vuelve el camino. Lo que estorba al camino se vuelve camino.''}\\[1mm]
Cada error de hoy, bien observado, se convierte en pieza del camino que pisarás mañana.
\end{softbox}

\small
\begin{tabularx}{\textwidth}{>{\bfseries}p{3.6cm}Y}
\rowcolor{milcGradeDark!12}Nombre completo: & \linead{10cm}\\
Fecha: & \linead{4cm} \quad Curso/grupo: \linead{4cm}\\
Mi firma: & \linead{9cm}\\
\end{tabularx}
\normalsize

\begin{infoband}{milcGradeDark}
\textbf{Para la próxima sesión.} Esta semana, antes de dormir, mira el reloj de tu cuaderno. Cumple los 3 compromisos que firmaste. La próxima clase vas a aprender sobre \textbf{privacidad y datos personales}: qué son, por qué importan, y cómo configurar tu celular y tus apps para protegerte. Hoy viste tu huella. La próxima ves cómo cerrar las puertas para que solo entre quien tú quieres.
\end{infoband}

\end{document}
